\scanpage{18}
% archive sheet number: 17
% @ text / head
Soit $\omega$ sur $S_{\unsure{\mathrm{gros}|\mathrm{zar}}}$, avec $F_\omega,V_\omega$, on lui associe

% @ display / upper third, item 1 and its right-hand addition
\[
1^\circ)\qquad M=\varphi^*(\omega),\qquad F_M=\varphi^*(F_\omega),\quad V_M=\varphi^*(V_\omega)
\]
\margin{$M=\omega^{\unsure{\flat|b|p}}$}

% @ text / upper third, item 2
$2^\circ)$ Triangle
\margin{$\varepsilon^*(\Delta(M))=\Delta^*(\omega)$ complexe can. ext. de $\omega^{\unsure{\sharp|\flat|\prime}}$ par $\omega^\flat$}

% @ diagram / item 2, upper centre
\diagram{\begin{tikzcd}[ampersand replacement=\&,column sep=small,row sep=large]
\& {\mathcal T_2^{(\omega)}=[\omega\mathrel{\overset{V_\omega}{\text{---}}}\omega^{\flat}]} \arrow[dl,"1"'] \& \\
{\mathcal T(\omega)\ \unsure{\not\approx|\not\simeq}\quad\mathcal T_0^{(\omega)}=[\omega^{V}\mathrel{\overset{F_\omega}{\longrightarrow}}\omega]}
\arrow[rr] \& \& {\mathcal T_1^{(\omega)}=[\omega^{(p)}\mathrel{\overset{0}{\text{---}}}\omega^{(p)}]} \arrow[ul]
\end{tikzcd}}

% @ text / below triangle
déduit de $\omega$ $\bigl(\mathcal T_1^{(\omega)}=\Delta^*(\omega)\bigr)$

% @ text / middle, item 3
$3^\circ)$ Isom \uuline{can}
% @ display / middle, alpha
\[
\varphi^*(\mathcal T(\omega))\overset{\alpha_\omega}{\simeq}\mathcal T(\varphi^*(\omega))
\]

% @ display / middle, item 4
\[
4^\circ)\quad\text{L'\uline{identité}}\qquad
\mathcal T(\omega)_1\overset{\beta_\omega}{\simeq}\varepsilon^*\Delta^*(M))\simeq
[\omega^{\unsure{\sharp|N}}\mathrel{\overset{0}{\text{---}}}\omega^\flat]
\]

% @ text / centre, above rule
\uline{Question} Le foncteur
% @ display / assignment
\[
(\omega,F_\omega,V_\omega)\longmapsto
(\varphi^*(\omega),\varphi^*(F_\omega),\varphi^*(V_\omega),\mathcal T(\omega),\alpha_\omega,\beta_\omega)
\]
% @ text / immediately above rule
est-il pl. fidèle ?? L'est-il du moins sur la sous-catégorie des triples $(\omega,F_\omega,V_\omega)$ avec $F_\omega=0$ ?
\par\noindent\rule{\linewidth}{0.4pt}

% @ diagram / below rule, left brace
\diagram{\begin{tikzcd}[ampersand replacement=\&,column sep=large,row sep=small]
{\left\{\vphantom{\begin{matrix}a\\a\\a\end{matrix}}\right.} \& \varphi^*(\omega)\arrow[r,"u"] \& \varphi^*(\omega') \& \text{compatible avec }F,V \\
\& \mathcal T(\omega)\arrow[r,"v"] \& \mathcal T(\omega') \& \text{hom de nos triangles}
\end{tikzcd}}

% @ text / below brace, two qualifications
$v_1$ déterminé comme $\varepsilon^*(u)$ \qquad rendant commutatifs les diagrammes

$\varphi^*(v)$ déterminé par $u$ comme $\mathcal T(u)$

% @ diagram / lower right, upper square
\diagram{\begin{tikzcd}[ampersand replacement=\&,row sep=small]
\mathcal T_1(\omega)\arrow[r,"v_1"]\arrow[d,Rightarrow,"{\beta_\omega\,\unsure{\wr}}"',"1"] \& \mathcal T_1(\omega')\arrow[d,Rightarrow,"\beta_{\omega'}"]\\
\varepsilon^*\Delta^*(\varphi^*(\omega))\arrow[d,equal] \& \varepsilon^*(\Delta^*(\varphi^*(\omega')))\arrow[d,equal]\\
\Delta^*(\omega)\arrow[r,"\varepsilon^*(u)"'] \& \Delta^*(\omega')
\end{tikzcd}}

% @ diagram / bottom right, lower square
\diagram{\begin{tikzcd}[ampersand replacement=\&]
\varphi^*(\mathcal T(\omega))\arrow[r,dash,"\alpha_\omega","\sim"']\arrow[d,"\varphi^*(v)"'] \& \mathcal T(\varphi^*(\omega))\arrow[d,"\mathcal T(u)"]\\
\varphi^*(\mathcal T(\omega'))\arrow[r,dash,"\alpha_{\omega'}","\sim"'] \& \mathcal T(\varphi^*(\omega'))
\end{tikzcd}}
